\clearpage
\section{Рефлексивный анализ выполненных заданий}

\subsection{Написание функции-расширения}

В первом задании была подготовлена ветка \texttt{task01}, создано решение \texttt{practice2025}, а также библиотека классов \texttt{task01} и тестовый проект \texttt{task01tests}. В статическом классе \texttt{StringExtensions} реализован метод расширения \texttt{IsPalindrome}. Для его вызова в форме метода экземпляра первый параметр объявлен как \texttt{this string input}. Обработка строки состоит из удаления пробельных символов и знаков препинания, приведения символов к единому регистру и сравнения нормализованной последовательности с её обратным представлением. После реализации были подготовлены тесты xUnit и настроена автоматическая сборка.

Основное затруднение было связано с определением поведения для граничных входных данных. Формального сравнения исходной строки с её перевёрнутой копией недостаточно, поскольку фразы-палиндромы содержат пробелы, знаки препинания и буквы разного регистра. Дополнительно требовалось определить результат для пустой строки, строки из пробелов и строки, состоящей только из пунктуации.

Проблема была решена посредством предварительной нормализации данных. Для классификации символов использованы \texttt{char.IsWhiteSpace} и \texttt{char.IsPunctuation}; после очистки проверяется наличие хотя бы одного значимого символа. Набор тестов был расширен граничными случаями. Такой подход позволил отделить подготовку входных данных от собственно проверки симметричности.

Наиболее сложным моментом оказалось формирование однозначного контракта метода для строк, которые после очистки становятся пустыми. Принятое решение возвращать \texttt{false} исключает формально истинный, но содержательно бессмысленный результат для пустой последовательности.

\subsection{Введение в LINQ}

Во втором задании создана ветка \texttt{task02}, проекты \texttt{task02} и \codebreak{task02tests}, а также ссылка тестового проекта на библиотеку. Модель \texttt{Student} содержит имя, факультет и список оценок. Класс \texttt{StudentService} реализует выбор студентов факультета, фильтрацию по среднему баллу, сортировку по имени, группировку по факультету и поиск факультета с максимальным средним баллом. Все операции с коллекциями выражены посредством LINQ.

Затруднение состояло в правильном выборе уровня агрегации. Средний балл отдельного студента вычисляется по его оценкам, тогда как средний балл факультета должен учитывать группу студентов. Ошибка в последовательности операций \texttt{GroupBy}, \texttt{Average} и \texttt{OrderByDescending} могла привести к сравнению несопоставимых величин. Отдельного внимания потребовало требование не использовать циклы.

Для решения задачи каждый метод был представлен как последовательность преобразований: фильтрация выполняется через \texttt{Where}, сортировка --- через \texttt{OrderBy}, группировка --- через \texttt{ToLookup}, а определение лучшего факультета --- через группировку, вычисление среднего значения и выбор первого элемента после сортировки по убыванию. Корректность каждого преобразования подтверждена отдельным тестом.

Наиболее сложным оказался метод поиска факультета с максимальным средним баллом, поскольку он объединяет группировку, вложенную агрегацию и выбор итогового ключа. Его декомпозиция в последовательность LINQ-операций сделала логику прозрачной и устранила необходимость в промежуточных циклах.

\subsection{Итераторы в C\#}

В третьем задании создана обобщённая коллекция \codebreak{CustomCollection<T>}, реализующая \texttt{IEnumerable<T>}. Внутреннее хранение организовано с помощью \texttt{List<T>}. Добавлены методы \texttt{Add} и \texttt{Remove}, стандартный перечислитель, обратный итератор \codebreak{GetReverseEnumerator}, генератор \codebreak{GenerateSequence}, а также LINQ-метод \texttt{FilterAndSort}. Проекты \texttt{task03} и \codebreak{task03tests} включены в общее решение.

Основное затруднение было связано с различием между готовой коллекцией и лениво вычисляемой последовательностью. Метод с \texttt{yield return} не формирует полный результат при вызове, а сохраняет состояние выполнения и выдаёт элементы по мере перечисления. Для обратного обхода требовалось корректно определить начальный индекс и условие завершения, особенно для пустой коллекции.

Обратный обход реализован от индекса \texttt{Count - 1} до нуля. Генератор числовой последовательности выдаёт ровно \texttt{count} значений, начиная с \texttt{start}. Фильтрация и сортировка делегированы стандартным операциям \texttt{Where} и \texttt{OrderBy}. Тесты материализуют ленивые последовательности методом \texttt{ToList}, после чего сравнивают фактический порядок элементов с ожидаемым.

Наиболее сложным моментом стало понимание отложенного характера выполнения итератора. Результат метода с \texttt{yield return} существует не как заранее заполненный список, а как правило последовательной выдачи значений. Это свойство особенно важно при работе с большими наборами данных и цепочками LINQ.

\subsection{Интерфейсы в C\#}

В четвёртом задании определён интерфейс \texttt{ISpaceship}, задающий операции движения, поворота и выстрела, а также свойства скорости и мощности огня. Классы \texttt{Cruiser} и \texttt{Fighter} реализуют единый контракт. Крейсер имеет скорость 50 и мощность 100, истребитель --- скорость 100 и мощность 50. Для проверки поведения добавлены свойства пройденного расстояния, текущего угла и количества выпущенных ракет с закрытыми сеттерами.

Затруднение состояло в том, что методы интерфейса имеют возвращаемый тип \texttt{void}. Простая запись сообщений в консоль позволила бы вызвать методы, однако не обеспечила бы надёжной проверки результата. Требовалось сохранить заданный интерфейс и одновременно сделать поведение наблюдаемым для модульных тестов.

Проблема решена за счёт инкапсулированного состояния конкретных классов. Метод \texttt{MoveForward} увеличивает расстояние на значение скорости, \texttt{Rotate} изменяет угол, а \texttt{Fire} увеличивает счётчик ракет. Закрытая запись свойств предотвращает произвольное изменение состояния извне. Тесты подтверждают как характеристики кораблей, так и изменения после команд.

Наиболее сложным оказалось согласование абстрактного контракта и проверяемой реализации. Интерфейс должен описывать возможности корабля, не раскрывая детали хранения состояния, тогда как тестам необходимо наблюдать последствия операций. Добавление доступных только для чтения извне свойств обеспечило требуемый баланс полиморфизма и инкапсуляции.

\subsection{Изучение класса с помощью рефлексии}

В пятом задании реализован класс \texttt{ClassAnalyzer}, принимающий объект \texttt{Type}. Средствами \texttt{System.Reflection} получаются публичные методы, имена параметров и возвращаемый тип выбранного метода, все поля, публичные свойства и сведения о наличии атрибута. Для выбора категорий членов применены флаги \texttt{BindingFlags}; обработка полученных массивов выполнена с помощью LINQ без циклов.

Главное затруднение было связано с точным выбором флагов поиска. Стандартные методы рефлексии могут включать унаследованные члены, пропускать статические элементы либо не возвращать закрытые поля. Кроме того, требование к \texttt{GetMethodParams} предусматривало не только имена параметров, но и возвращаемый тип.

Для методов использована комбинация \texttt{Public}, \texttt{Instance}, \texttt{Static} и \texttt{DeclaredOnly}. Для полей дополнительно указан \texttt{NonPublic}. Имена параметров формируются оператором \texttt{Select}, а имя возвращаемого типа присоединяется оператором \texttt{Append}. Если метод отсутствует, последовательность остаётся пустой. Атрибут определяется с помощью обобщённого метода \texttt{GetCustomAttribute<T>}.

Наиболее сложным моментом стало понимание того, что результат рефлексии зависит не только от анализируемого типа, но и от заданной области поиска. Явное указание \texttt{BindingFlags} сделало поведение анализатора предсказуемым и позволило проверить публичные, приватные, статические и экземплярные члены отдельными утверждениями.
